\section{元朝的尺度：行省、江南与多语帝国}

忽必烈建元、定都大都和南宋覆亡，使一个跨越草原、农耕区、绿洲与海域的帝国获得了新的中国王朝名称；这不意味着统治从此只有一个中心或一种社会。行中书省、路府、站赤、宣慰与地方军政机构，都是把命令、税粮、案件和人员在不同区域间接起来的办法。它们的边界、权限与实际覆盖会因战事、运输、地方势力和中枢调整而变动。因此，行省不宜被写成现代省界的预演，更应被理解为征服与治理之间不断校准的行政场域。

\subsection{江南不是被动的财政仓库}

元初接收江南，需要处理的是南宋留下的田赋、仓储、盐业、港口、市镇、士人和地方组织。江南税粮可以支持大都与北方供给，但粮食要经过征收、仓储、河运、海运和交接才能成为国家资源；每一个环节都可能因水位、盗抢、折耗、地方议价或官员舞弊而改变。把江南简称为“财赋基地”，能指出其重要性，却容易掩盖这种运输政治及其对地方家庭的压力。

纸钞同样说明制度的双面性。它为跨地域支付和政府调度提供便利，也依赖发行、兑换、物价与信任。财政紧张时增加发行或改变折纳方式，可能把中枢的困难转移给市场和纳税者。盐课、商税、包银、漕粮和劳役不是单独的部门问题，而是人们在国家账目、地方市场和家庭生计之间反复折算的关系。不同地区留下的账簿和文书远不均衡，故不能用一个数字概括整个帝国的繁荣或衰退。

\subsection{身份、宗教与流动}

元代的身份分类和任官秩序确实塑造了政治机会，却不能把“蒙古、色目、汉人、南人”等行政类别直接等同于固定的生活经验。军人、译者、工匠、商人、僧侣、儒士、妇女与寄居者，在籍贯、职业、法律案件和宗教网络中形成了交错关系。多语碑刻、契约与寺院材料使其中一部分人可见，也提醒我们：不会留下文书的人并不因此没有改变制度。

佛教、道教、伊斯兰教和地方信仰的机构既与国家资源相连，也可能成为土地、劳力、慈善、教育和跨区交往的节点。朝廷的保护、限制或征调会改变这些机构的地位；把宗教只写成皇帝的信仰，或只写成社会的避难所，都会错过它们实际管理财产和人群的能力。海港和陆路贸易亦然：帝国范围提供机会，航海风险、地方管制和战争又不断限定机会。

\subsection{交读窗口：杂剧法庭不能当作案卷}

关汉卿《感天动地窦娥冤》约成于十三世纪后期，具体写作年月与演出场合都无从确证。作为元杂剧，它用唱词、科介和高度安排的冤狱推进冲突；“地也，你不分好歹何为地！天也，你错勘贤愚枉做天！”把法律失效写成天地倒置。它的历史照明力正在于展示元初汉文舞台可以怎样想象贫困妇女、媒介人、官府与申冤，而不是提供一宗可供统计的真实判案。

《元史》、《元典章》和《通制条格》能够分别给出机构、条格与规范语言；地方契约、碑刻、寺院文书和偶存判牍才可局部检验诉讼、担保与财产关系实际如何运作。八思巴字、蒙古文、畏兀儿文与波斯文材料能打破只从汉文舞台观察帝国的局限，但存量、译名、地域和保存条件极不平衡，也很少能同一出戏的情节逐项相证。钱穆关于制度是否有稳定通道的追问，在这里须同吕思勉对家庭、职业与弱势者处境的追问相接；《窦娥冤》只证明一套戏剧化的道德想象，不能证明元代司法必然如此，或所有女性都拥有相同的申诉语言。

\subsection{河工、灾害与终局}

十四世纪的河患、漕运困难、财政信用与政治竞争彼此交织。一三五一年河工动员与红巾军起事同年发生，是理解危机的重要并置，不能因此断言一项工程单独引发了所有叛乱。各地起事者、地方武装与元廷官军的目标并不相同；疾病、歉收、赋役和交通断裂又会在不同地区产生不同后果。到一三六八年元廷退出大都，王朝在中原的统治结束，北方与欧亚网络中的蒙古政治却没有随之消散。

《元史》是在明初仓促修成的官修史，擅长保存帝纪和部分制度线索，也带有新王朝整理旧政权的框架。《元典章》《通制条格》、碑刻、契约、地方志、波斯文旅行与行政材料能补充不同层面的记录，却各有地域、语言和保存偏差。元史研究最重要的收益之一，正是把这些材料放在一起，既不把行省与纸钞神化为全能制度，也不把帝国的困难简化成某个民族或某位皇帝的性格。
